\section{What can overlapping patches predict?}\label{sec:patches}

A bit that is zero with probability \(0.9\) can be predicted with \(90\%\)
accuracy without reading any helper. The information supplied by an allowed
helper set \(B\) is \(I(X_q;X_B)=H(X_q)-H(X_q\mid X_B)\). In the prediction credit
\[
  1-H(X_q\mid X_B)=[1-H(X_q)]+I(X_q;X_B),
\]
the first term is due to the marginal bias and the second measures the
gain from observing the helpers. The geometry of the relation bounds
the sum of these gains, even when different
patches use the same bits. Write \(X_S=(X_v)_{v\in S}\) and use entropy
\(H\) in bits. The Shannon entropy in the Lean declarations is measured in
nats, so its value in bits is obtained by division by \(\ln 2\).

Let \(V\) be a finite set of coordinates. Each member of a finite indexed family of \emph{patches} \(P_j\subseteq V\)
has a set \(A_j\subseteq P_j\) of scored targets. For each \(q\in A_j\),
choose the helpers \(B_{jq}\subseteq P_j\) available to its predictor.
A target scored in two patches contributes two occurrences \((j,q)\).
Form the bipartite graph on the two copies of \(A_j\) by joining \(L_v\) to
\(R_q\) when \(v\in B_{jq}\cap A_j\). A patch satisfies the
\emph{separation condition} if \(L_q\) and \(R_q\) lie in different
components of this graph for every \(q\in A_j\).

An ambient distance bound gives a sufficient separation test. Take a
relation \(N\). Suppose \(B_{jq}\subseteq N(q)\cap P_j\),
\(\card{A_j}\le w\) with \(w\ge1\), and
\(\dist(L_q,R_q)>2w-1\) for every scored occurrence. A connection inside
the patch would contain a simple path on at most \(2\card{A_j}\) ports.
The same path in the ambient graph would have length at most \(2w-1\).
Thus the separation condition holds. The patch can contain arbitrarily
many unscored helpers.

\begin{theorem}[Theorem P, information supplied by overlapping patches]
\label{thm:patch-credit}
\formal{TapeStorageRates/PatchResidual.lean}{TapeStorageRates.weighted_patch_capacity,TapeStorageRates.patch_capacity_of_long_targets}
Let \(h\) be a polymatroid rank on \(V\), that is, a normalized, monotone,
submodular set function. Write \(h(q\mid B)=h(B\cup\{q\})-h(B)\).
Give each coordinate a capacity \(c_q\ge h(\{q\})\) and each patch a
weight \(\lambda_j\ge0\). Suppose the separation condition holds, and
\(\sum_{j:v\in P_j}\lambda_j\le r\) for every coordinate \(v\). Then
\[
 \sum_j\lambda_j\sum_{q\in A_j}\bigl(c_q-h(q\mid B_{jq})\bigr)
 \le r\left[\left(\sum_{v\in V}c_v-h(V)\right)
              +\left(\sum_{v\in V}h(\{v\})-h(V)\right)\right].
\]
For a finite random vector \(X\), write
\(\mathrm{TC}(X)=\sum_v H(X_v)-H(X_V)\) for its total correlation.
Taking \(h(S)=H(X_S)\) and \(c_v=H(X_v)\) gives
\[
 \sum_j\lambda_j\sum_{q\in A_j} I(X_q;X_{B_{jq}})
 \le 2r\,\mathrm{TC}(X).
\]
\end{theorem}
\begin{outline}
Fix one patch and write \(P=P_j\), \(A=A_j\), \(D=P\setminus A\).
Condition on \(D\). All remaining helpers of a target \(q\) have their
left ports in the component of \(R_q\). Order the components and first
supply the targets whose left component precedes their right component.
Then process the other targets in order of their left components. For
each such target its right component came earlier, so all of its helpers
are available. A supplied target costs at most \(h(q\mid D)\) in the
chain rule, and a target reached with its helpers costs at most
\(h(q\mid B_{jq})\). Reversing the order exchanges the two roles for
each target. Adding the bounds gives
\[
 2\bigl(h(P)-h(D)\bigr)
 \le \sum_{q\in A} h(q\mid D)+\sum_{q\in A}h(q\mid B_{jq}).
\]
The formal proof groups the targets by their left components to obtain
a single elimination piece in the sense of~\cite{PaperSeven}. The
common-observations corollary of the soft-recovery memo gives this
inequality~\cite[Corollary~3.2]{SoftRecovery}. The two component orders explain the
same factor two as in the exact reconstruction argument
of~\cite[Lemma~4.1]{PaperFive}.

Use \(h(q\mid D)\le h(\{q\})\) and bound \(h(D)\) both by
\(\sum_{v\in D} h(\{v\})\) and by \(\sum_{v\in D} c_v\).
The contribution of this patch is at most the sum of the two deficits
on the right side of the theorem, evaluated on \(P\).
To sum either deficit, let \(a_v\) stand for \(c_v\) or \(h(\{v\})\)
and fix an order on \(V\), with \(V_{<v}\) its initial segment before
\(v\). Diminishing returns gives
\[
 \sum_{v\in P}a_v-h(P)\le\sum_{v\in P}
                   \bigl(a_v-h(v\mid V_{<v})\bigr).
\]
On weighting and summing, each nonnegative term on the right has
coefficient at most \(r\). The chain rule on \(V\) gives the claim.
The formal summation in \lean{TapeStorageRates.patch_deficit_packing}
uses the weighted submodular cover inequality from the Lean library
of~\cite{PaperTwo}.
\end{outline}

\begin{corollary}[Errors of local predictors]\label{cor:patch-prediction}
\formal{TapeStorageRates/PatchPrediction.lean}{TapeStorageRates.PatchPredictor.prediction_capacity,TapeStorageRates.PatchPredictor.prediction_deficit}
Let \(X\) be a random binary word on \(n=\card V\) coordinates.
Use unit patch weights and the ambient distance test above. Let \(r\)
bound the number of patches containing any coordinate. Each occurrence
\(i=(j,q)\) uses a fixed Boolean predictor of \(X_q\) from \(X_{B_{jq}}\)
with error probability \(e_i\). Assume the number \(m\) of scored occurrences is positive. Set
\(\bar e=m^{-1}\sum_i e_i\), and \(\delta=n-H(X_V)\).
With \(h_2(t)=-t\log_2t-(1-t)\log_2(1-t)\) and \(0\log_2 0=0\),
\[
 m[1-h_2(\bar e)]
 \le \sum_i[1-h_2(e_i)]
 \le r\bigl(\delta+\mathrm{TC}(X)\bigr)
 \le 2r\delta .
\]
\end{corollary}
\begin{outline}
The predictor and its error bit together determine \(X_q\). Thus
\(H(X_q\mid X_{B_{jq}})\le h_2(e_i)\), so \(c_q-h_2(e_i)\)
is a lower bound on the contribution of occurrence \(i\) in
\cref{thm:patch-credit} for every choice of the capacities. Use \(c_v=1\) in
\cref{thm:patch-credit} and \(H(X_v)\le1\). Concavity of binary
entropy gives the bound in terms of \(\bar e\).
\end{outline}

Equality holds at every error level \(0\le e\le1\).
In one patch, score both coordinates of \((U,U\mathbin{\oplus}Z)\),
where \(U\) is a fair bit and \(Z\) is an independent bit with
\(\Pr(Z=1)=e\). Each predictor copies its partner. Both paired
distances are infinite, \(r=1\), both errors equal \(e\), and the joint
entropy is \(1+h_2(e)\). Thus the two sides of the bound
\(m[1-h_2(\bar e)]\le2r\delta\) equal \(2[1-h_2(e)]\).
For \(e\ne 1/2\) this proves that the factor two is necessary for general binary laws. The example is \lean{TapeStorageRates.noisy_pair_prediction_sharp}.

For a uniform word in a binary code of length \(n>0\), suppose
\(\delta/n\le D/s\) with \(D,s>0\). If \(m\ge\alpha n\), then
\[
 r\ \ge\ \frac{\alpha[1-h_2(\bar e)]}{2D}\,s.
\]
This is \lean{TapeStorageRates.PatchPredictor.overlap_of_near_unit_rate}.
Predicting a positive density of separated targets with a fixed advantage
over a fair guess therefore requires overlap growing at least linearly
in \(s\). For the families constructed in \cref{thm:main}, take the radius
\(2w-1\) after fixing \(s\) and \(w\). For a native fibre supplied by
\cref{thm:native-tail}, \(X\) can be taken uniform on that particular
fibre, using its own actual permitted helpers and entropy deficit.

\begin{corollary}[Exact recovery over a finite alphabet]
\label{prop:patch-coverage}
\formal{TapeStorageRates/Patches.lean}{TapeStorageRates.recovery_patch_coverage}
Let \(C\subseteq\Sigma^V\) be nonempty, where \(a=\card{\Sigma}\ge2\).
Each patch \(P_j\) serves at most \(w\ge1\) targets in \(A_j\subseteq P_j\).
For each \(q\in A_j\), its value on \(C\) is determined by the
coordinates in \(N(q)\cap P_j\). Each coordinate belongs to at most
\(r\) patches. Let \(S\subseteq\bigcup_j A_j\), and let \(b_N(t)\)
count all coordinates in \(V\) at paired distance at most \(t\).
Then
\[
 \card S\le b_N(2w-1)+2r\bigl(n-\log_a\card C\bigr).
\]
In fact the last term alone bounds the number of distinct targets in
\(S\) at distance greater than \(2w-1\).
\end{corollary}
\begin{outline}
Give \(C\) the uniform law. In each patch score only the served targets
at distance greater than \(2w-1\). Their conditional entropies given the
permitted helpers vanish. In \cref{thm:patch-credit} take
\(c_v=\log_2 a\). Each long occurrence contributes \(\log_2 a\), and
every distinct long target in \(S\) contributes at least once. Both
global deficits are at most \(n\log_2 a-\log_2\card C\).
\end{outline}

Thus even if the alphabet varies with the code, a family with a
vanishing fraction of short pairs at radius \(2w-1\) has upper
limiting covered fraction at most \(2r(1-\liminf_i R_i)\), where
\(R_i=\log_{\card{\Sigma_i}}\card{C_i}/\card{V_i}\) and \(w,r\)
are fixed. This is \lean{TapeStorageRates.patch_coverage_limsup}.
Disjoint repeated pairs attain the finite bound at \(r=1\), as in
\lean{TapeStorageRates.repetition_pairs_sharp}.

\begin{corollary}[Input-dependent choice of predictors]\label{cor:patch-menu}
\formal{TapeStorageRates/PatchPredictionEvent.lean}{TapeStorageRates.patch_prediction_menu_probability,TapeStorageRates.patch_prediction_event_probability}
Let \(E\subseteq\{0,1\}^V\), \(n=\card V\), \(0\le e\le1/2\), and
\(r\ge1\). Fix a menu of at most \(2^h\) predictor systems satisfying
the ambient distance test with common bounds \(w,r\). A system specifies
its relation, patches, scored targets, helpers and prediction rules.
Suppose that, for each \(x\in E\), some system scores at least \(m\)
occurrences and errs on at most a fraction \(e\) of all the occurrences
it scores on \(x\). Then
\[
 \frac{\card E}{2^n}\le
 2^{\,h-m[1-h_2(e)]/(2r)}.
\]
In particular a transcript selecting the system can be charged
\(h=\log_2\card{\{\text{labels attained on }E\}}\) when \(E\ne\varnothing\).
\end{corollary}
\begin{outline}
For each system, restrict to the inputs of \(E\) on which it meets
the two requirements. Under the uniform law on any such nonempty subset,
its mean error is at most \(e\). By \cref{cor:patch-prediction} and
monotonicity of \(h_2\) on \([0,1/2]\), the subset has at most
\(2^{\,n-m[1-h_2(e)]/(2r)}\) elements. Their union contains \(E\).
\end{outline}

A tape lower bound along this route still requires a proof that
successful executions admit predictors with the stated accuracy and
overlap. On a fixed native fibre one can test the predictors under its
uniform law. On the full cube one must also bound the number of distinct
systems selected by the trace, the genie word and any other observation
used to choose the targets.
